\section*{Chapitre \ref{ch:b3:bacteriology} --- Bactériologie~: croissance, physiologie et génétique}

\begin{solution}{exo:b3:bacteriology:1}
\omterm{def:b3:bacteriology:envelope}{Gram positif}~: membrane cytoplasmique enveloppée d'une paroi de
\omterm{def:b3:bacteriology:envelope}{peptidoglycane} épaisse et multicouche, traversée d'acides téichoïques.
\omterm{def:b3:bacteriology:envelope}{Gram négatif}~: une couche mince de \omterm{def:b3:bacteriology:envelope}{peptidoglycane} dans un périplasme,
puis une \omterm{def:b3:bacteriology:envelope}{membrane externe} dont le feuillet externe est du
\omterm{def:b3:bacteriology:envelope}{lipopolysaccharide}, avec des porines. Le complexe violet de
gentiane--iode est retenu dans la paroi épaisse et déshydratée par
l'alcool du \omterm{def:b3:bacteriology:envelope}{Gram positif} (violet) mais lessivé hors du \omterm{def:b3:bacteriology:envelope}{Gram négatif} au
travers de sa paroi mince et de sa \omterm{def:b3:bacteriology:envelope}{membrane externe} désorganisée, qui
prend alors la contre-coloration rose.
\end{solution}

\begin{solution}{exo:b3:bacteriology:2}
Un facteur $10^{6}$ $= 2^{19.9}$~: environ $20$ générations~; $T = \qty{8}{h}/20
= \qty{24}{min}$~; $\mu = \ln 2/\qty{0.4}{h} = \qty{1.7}{h^{-1}}$.
\end{solution}

\begin{solution}{exo:b3:bacteriology:3}
\omterm{def:b3:bacteriology:hgt}{Transformation} --- Griffith (1928) puis Avery, MacLeod et McCarty
(1944)~: des pneumocoques virulents tués par la chaleur transforment des
pneumocoques vivants avirulents, et l'agent est de l'ADN. \omterm{def:b3:bacteriology:hgt}{Transduction}
--- Zinder et Lederberg (1952)~: le phage P22 porte des gènes de
\emph{Salmonella} d'une souche à l'autre au travers d'un filtre qui
arrête les cellules. \omterm{def:b3:bacteriology:hgt}{Conjugaison} --- Lederberg et Tatum (1946)~: deux
souches auxotrophes d'\emph{E.~coli} ne donnent de recombinants
prototrophes qu'une fois mélangées, le contact entre cellules étant
requis (le tube en U de Davis).
\end{solution}

\begin{solution}{exo:b3:bacteriology:4}
Pénicilline~: les transpeptidases qui réticulent le \omterm{def:b3:bacteriology:envelope}{peptidoglycane}, que
les cellules animales n'ont pas. Streptomycine~: la petite sous-unité
ribosomique, dont la forme bactérienne diffère de l'eucaryote.
Ciprofloxacine~: l'ADN gyrase, topo-isomérase bactérienne absente chez
l'homme. Rifampicine~: la sous-unité $\beta$ de l'ARN polymérase
bactérienne, que l'enzyme eucaryote ne partage pas.
\end{solution}

\begin{solution}{exo:b3:bacteriology:5}
$S^{*} = K_{S}D/(\mu_{\max} - D)$, $X^{*} = Y(S_{0} - S^{*})$. $D = 0.2$~:
$S^{*} = \qty{0.0067}{g/L}$, $X^{*} = \qty{0.80}{g/L}$. $D = 0.6$~:
$S^{*} = 0.06$, $X^{*} = 0.78$. $D = 0.79$~: $S^{*} = 1.58$, $X^{*} = 0.17$.
Lessivage quand $D = \mu(S_{0}) = 0.8\times 2/2.02 = \qty{0.79}{h^{-1}}$.
\end{solution}

\begin{solution}{exo:b3:bacteriology:6}
$(N/V)_{\text{seuil}} = kA_{\text{seuil}}/p$. Organe lumineux~: $0.01\times
6\times 10^{12}/100 = 6\times 10^{8}$ par litre $= 6\times 10^{5}$ par
mL. Eau de mer~: $10\times 6\times 10^{12}/100 = 6\times 10^{11}$ par
litre $= 6\times 10^{8}$ par mL --- mille fois plus, jamais atteint en
mer.
\end{solution}

\begin{solution}{exo:b3:bacteriology:7}
\qty{25}{mm}~: sensible, CMI basse. \qty{12}{mm}~: sensibilité réduite,
intermédiaire --- peut échouer aux doses ordinaires. \qty{0}{mm}~:
résistante, croissance jusqu'au disque. Le médicament diffuse depuis le
disque de sorte que sa concentration décroît avec la distance de façon
reproductible~; la croissance s'arrête là où la concentration égale la
CMI, si bien que le rayon de la zone et la CMI sont liés par une courbe
d'étalonnage établie avec des souches de CMI connue.
\end{solution}

\begin{solution}{exo:b3:bacteriology:8}
Le \omterm{def:b3:genetic-engineering:tools}{plasmide} conjugatif porte un intégron dont le jeu de cassettes
contient plusieurs gènes de résistance, plus des transposons qui en
portent d'autres~; le \omterm{def:b3:genetic-engineering:tools}{plasmide} passe en un seul accouplement et les
exprime tous. Ils se groupent parce que la sélection par l'un
quelconque des médicaments maintient le \omterm{def:b3:genetic-engineering:tools}{plasmide} entier, de sorte que
les gènes qui y font du stop persistent (cosélection)~; les intégrons
capturent des cassettes~; les transposons sautent sur les \omterm{def:b3:genetic-engineering:tools}{plasmides}~; et
c'est le \omterm{def:b3:genetic-engineering:tools}{plasmide}, non le gène, qui est l'unité de transfert.
\end{solution}

\begin{solution}{exo:b3:bacteriology:9}
Le médicament seul~: $10^{-7}\times 10^{9} = 100$ cellules résistantes
attendues. Deux médicaments indépendants~:
$10^{-14}\times 10^{9} = 10^{-5}$. Les persistantes ne sont pas des
mutantes~: une petite
fraction de cellules dormantes survit à n'importe quel médicament et
repeuple une population entièrement sensible une fois le médicament
parti, de sorte qu'un second médicament n'y change rien~; seul un
traitement assez long pour les attraper à leur réveil, ou un médicament
actif sur les cellules dormantes, y parvient.
\end{solution}

\begin{solution}{exo:b3:bacteriology:10}
$X > X^{*}$~: la consommation $\mu X/Y$ dépasse l'apport, donc $S$
tombe sous $S^{*}$~; alors $\mu(S) < D$ et
$\mathrm{d}X/\mathrm{d}t < 0$, $X$ redescend~; à mesure que $X$ baisse,
$S$ remonte. Les deux écarts sont corrigés~: l'état stationnaire est
stable. Deux espèces~: chacune a besoin de $\mu_{i}(S) = D$ pour
persister, ce qui fixe son propre $S_{i}^{*}$~; le nutriment ne peut
prendre qu'une valeur, donc une espèce au plus persiste. Celle dont le
$S^{*}(D)$ est le plus bas fait tomber $S$ sous ce dont l'autre a
besoin, et l'autre est lessivée~; laquelle c'est peut changer avec $D$
si leurs courbes de Monod se croisent.
\end{solution}

\begin{solution}{exo:b3:bacteriology:11}
Avec la densité $\rho = N/V$, l'état éteint $A = p_{0}\rho/k$ existe
tant qu'il reste sous $A_{\text{seuil}}$, c'est-à-dire
$\rho < kA_{\text{seuil}}/p_{0}$~; l'état allumé $A = p_{1}\rho/k$ existe tant
qu'il reste au-dessus, c'est-à-dire $\rho
> kA_{\text{seuil}}/p_{1}$. Pour $kA_{\text{seuil}}/p_{1} < \rho <
kA_{\text{seuil}}/p_{0}$ les deux sont cohérents~: une population qui
s'est allumée se maintient allumée par sa production plus forte, une
population qui ne l'a pas fait reste éteinte. Usage~: l'engagement ---
une fois qu'une colonie a décidé de bâtir un biofilm ou de libérer une
toxine, un creux de densité ne défait pas la décision, et
l'interrupteur est net et synchrone plutôt que progressif.
\end{solution}

\begin{solution}{exo:b3:bacteriology:12}
(a) Fortes doses en cures courtes~: la concentration reste au-dessus de
la CMI des mutants et tue également le type sauvage et les mutants de
premier rang --- à retenir, dans la limite de la toxicité. (b) Faibles
doses en cures longues~: la concentration séjourne des jours dans la
fenêtre et fabrique de la résistance --- contre. (c) Arrêter quand les
symptômes cèdent~: les survivantes sont les cellules les moins
sensibles et la concentration qui redescend traverse la fenêtre ---
contre. (d) Association~: une cellule doit résister aux deux, ce qu'un
milliard de cellules ne contient pas --- pour.
\end{solution}

\begin{solution}{pb:b3:bacteriology:1}
\textbf{1.} $18$ générations~: $2^{18} = 2.6\times 10^{5}$ cellules,
$2.6\times 10^{3}$ par mL.
\textbf{2.} $2\times 10^{9}\times 100 = 2\times 10^{11}$ cellules $=
2^{37.5}$~: au bout de \qty{12.5}{h}~; masse \qty{0.2}{g}.
\textbf{3.} $6\times 10^{27}/10^{-12} = 6\times 10^{39} = 2^{132}$
cellules~: $132$ générations, \qty{44}{h}. Les nutriments, le dioxygène
et les déchets l'arrêtent en moins d'un jour~; la croissance
exponentielle est toujours brève.
\textbf{4.} Les cellules stationnaires ont démonté leurs ribosomes et
réprimé leurs enzymes métaboliques~; il leur faut les rebâtir avant de
se diviser. Les cellules exponentielles arrivent tout équipées.
\textbf{5.} $\mu = \ln 2/\qty{0.333}{h} = \qty{2.08}{h^{-1}}$~; à $T =
\qty{60}{min}$, \qty{0.69}{h^{-1}}~; après \qty{6}{h} $2^{6} = 64$ au
lieu de $2^{18}$~: un facteur $2^{12} \approx 4000$ de moins.
\textbf{6.} $200/(0.1\times 10^{-6}) = 2\times 10^{9}$ cellules viables
par mL. Le microscope compte aussi les cellules mortes et celles qui ne
pousseront pas sur la boîte, et un amas ne donne qu'une colonie.
\textbf{7.} $D = 0.5/1 = \qty{0.5}{h^{-1}}$~; $T = \ln 2/0.5 =
\qty{1.39}{h} = \qty{83}{min}$.
\textbf{8.} $S^{*} = 0.01\times 0.5/0.5 = \qty{0.01}{g/L}$~; $X^{*} =
0.5\times 1.99 \approx \qty{1}{g/L} = 10^{12}$ cellules par litre, soit
$10^{9}$ par mL.
\textbf{9.} $0.5\times 10^{12} = 5\times 10^{11}$ cellules par heure~;
$720/1.39
\approx 520$ générations par mois.
\textbf{10.} $D = 0.95$~: $S^{*} = 0.01\times 0.95/0.05 = \qty{0.19}{g/L}$,
$X^{*} = 0.5\times 1.81 = \qty{0.9}{g/L}$. À $F = \qty{1.1}{L/h}$, $D
> \mu_{\max}$~: lessivage.
\textbf{11.} Le mutant se contente de $S^{*} = 0.01\times 0.5/0.7 = \qty{0.0071}{g/L}$~;
à ce glucose la souche d'origine croît à $1.0\times 0.0071/0.0171 =
\qty{0.42}{h^{-1}} < D$ et décline comme $e^{-0.08t}$~: lessivée en
quelques semaines. Le mutant prend la place.
\textbf{12.} $10^{-9}\times 10^{12} = 1000$ mutations avantageuses par
génération~: l'adaptation est continue, un nouveau clone avantageux
balayant la cuve toutes les quelques centaines de générations --- le
chémostat est une machine à évoluer.
\textbf{13.} $kA_{\text{seuil}}/p = 0.05\times 6\times 10^{12}/500 =
6\times 10^{8}$ par litre $= 6\times 10^{5}$ cellules par mL.
\textbf{14.} $10^{10}/6\times 10^{5} \approx 1.7\times 10^{4}$ fois.
\textbf{15.} $20\times 6\times 10^{12}/500 = 2.4\times 10^{11}$ par
litre, $2.4\times 10^{8}$ par mL --- six ordres de grandeur au-dessus
des $10^{2}$ par mL de la mer~: sombres.
\textbf{16.} Un facteur dix fait $3.3$ doublements, \qty{6.6}{h}. À
$10^{9}$ par mL la population est encore $1700$ fois au-dessus du
seuil~: la lumière reste allumée.
\textbf{17.} En mer personne ne voit la lumière et elle ne nourrit aucun
hôte~: un mutant sombre économise \qty{20}{\%} et dépasse les
lumineuses. Dans l'organe, le calmar sélectionne la lumière --- il
expulse les tricheuses sombres ou cesse de les nourrir --- de sorte que
le coût achète un logis.
\textbf{18.} Bloquer le récepteur ou détruire l'\omterm{def:b3:bacteriology:quorum}{auto-inducteur}
(l'extinction du quorum)~: les bactéries vivent mais ne déploient
jamais leurs toxines, et le système immunitaire les élimine. Puisque
bactéries sensibles et «~résistantes~» croissent également --- rien
n'est tué --- l'avantage sélectif d'échapper au médicament est faible,
et la résistance devrait se répandre plus lentement que face à un
médicament qui tue.
\textbf{19.} A~: $10^{-8}\times 10^{9} = 10$~; B~: $100$~; les deux~:
$10^{-15}\times
10^{9} = 10^{-6}$.
\textbf{20.} $P_{0}(\text{A}) = e^{-10} = 5\times 10^{-5}$~;
$P_{0}(\text{les deux}) = e^{-10^{-6}} = 0.999999$.
\textbf{21.} $16 \to 8$~: une demi-vie, \qty{4}{h}~; $16 \to 1$~:
quatre, \qty{16}{h}~; dans la fenêtre de \qty{4}{h} à \qty{16}{h}, soit
\qty{12}{h} par dose.
\textbf{22.} Une dose toutes les \qty{12}{h}~: dans la fenêtre de
l'heure $4$ à l'heure $12$, les deux tiers du temps~; sur dix jours
les mutants de premier rang de A --- une dizaine au départ --- sont
sélectionnés et A seul échoue.
\textbf{23.} Au-dessus de \qty{8}{mg/L} en permanence~: une dose à
chaque demi-vie (\qty{4}{h}), ou une dose atteignant un pic de
\qty{64}{mg/L} toutes les \qty{12}{h} ($64 \to 8$ en trois demi-vies),
ou une perfusion continue. Coûts~: la toxicité de pics élevés, ou six
prises par jour que les patients ne tiendront pas.
\textbf{24.} $10^{-4}\times 10^{9} = 10^{5}$ persistantes survivent quoi
qu'il arrive~; à l'arrêt du traitement elles se réveillent et
reconstituent une population entièrement sensible --- une rechute. Le
traitement doit durer assez pour les attraper à leur sortie de
dormance, ce qui explique que la tuberculose demande six mois.
\textbf{25.} $10^{9}$ cellules par mL dans le chémostat~; la lumière
s'allume à $6\times 10^{5}$ cellules par mL~; $P_{0} = 0.999999$ que la
cure à deux médicaments ne rencontre aucune cellule résistante.
\end{solution}
